% 钱德拉塞卡拉·拉曼（综述）
% license CCBYSA3
% type Wiki

本文根据 CC-BY-SA 协议转载翻译自维基百科\href{https://en.wikipedia.org/wiki/C._V._Raman}{相关文章}

钱德拉塞卡拉·文卡塔·“C. V.”·拉曼爵士（Sir Chandrasekhara Venkata "C. V." Raman，/ˈrɑːmən/，\(^\text{[2]}\)音译“拉曼”；泰米尔语：சந்திரசேகர வெங்கட ராமன்，罗马化：Cantiracēkara Veṅkaṭa Rāmaṉ；1888年11月7日—1970年11月21日）是一位印度物理学家\(^\text{[3]}\)，以其在光散射领域的研究而闻名。\(^\text{[4]}\)他使用自己研制的分光仪，与学生K. S. 克里希南共同发现：当光线穿过透明物质时，被偏转的光会改变其波长。这种此前未知的光散射现象，他们称之为“修饰散射”，后来被命名为“拉曼效应”或“拉曼散射”。1930年，拉曼因这一发现获得诺贝尔物理学奖，成为第一位获得该奖项的亚洲人和非白人。\(^\text{[5]}\)

拉曼出生于泰米尔婆罗门家庭，自幼聪颖过人。他分别在11岁和13岁时，就完成了圣阿洛伊修斯英印高级中学的中学和高中课程。16岁时，他以优异成绩从马德拉斯大学附属的总督学院毕业，并以物理学荣誉学士的身份名列榜首。他的第一篇研究论文《光的衍射》发表于1906年，当时他仍是一名研究生。次年，他获得了硕士学位。19岁时，拉曼进入加尔各答的印度财政服务部门，担任副会计长。在那里，他结识了印度科学促进会（Indian Association for the Cultivation of Science，简称 IACS）——印度首个科研机构。该机构允许他独立开展研究，也正是在这里，拉曼在声学与光学领域作出了重要贡献。

1917年，拉曼受阿修托什·穆克吉任命，出任加尔各答大学下属拉贾巴扎尔理学院的首任帕利特物理学教授。在他首次前往欧洲的旅途中，亲眼目睹地中海的蓝色使他产生了强烈的好奇心——他认为当时流行的解释（即海洋的蓝色是由于天空中瑞利散射光的反射所致）是错误的。1926年，他创办了《印度物理学杂志》。1933年，他迁往班加罗尔，成为印度科学研究所的首位印度籍所长。同年，他还创立了印度科学院。1948年，他建立了拉曼研究所，并在此工作直至生命的最后时光。

拉曼效应于1928年2月28日被发现。印度政府将这一天定为“全国科学日”，每年举行纪念活动。
\subsection{早年生活与教育}
拉曼出生于英属印度马德拉斯省的蒂鲁吉拉帕利（Tiruchirappalli，今印度泰米尔纳德邦的蒂鲁吉拉帕利），父母均为泰米尔·艾耶尔婆罗门——父亲钱德拉塞卡拉·拉曼纳森·艾耶，\(^\text{[6][7]}\)母亲帕尔瓦蒂·阿玛尔。\(^\text{[8]}\)他是家中八个孩子中的第二个。\(^\text{[9]}\)他的父亲是一名中学教师，收入微薄。拉曼曾幽默地回忆说：“我出生时嘴里含着一把铜勺。因为在我出生时，父亲那‘辉煌’的薪水是每月十卢比！”\(^\text{[10]}\)1892年，全家迁往安得拉邦的维萨卡帕特南（当时称为Vizagapatam或Vizag），因为他的父亲被任命为A.V. 纳拉辛哈·饶夫人学院物理系教师。\(^\text{[11]}\)

拉曼在维萨卡帕特南的圣阿洛伊修斯英印高级中学接受教育。\(^\text{[12]}\)他在11岁时通过了中学毕业考试，13岁时以奖学金身份通过了艺术初级考试（First Examination in Arts，相当于今日的大学预科课程），\(^\text{[9][13]}\)并在两项考试中均名列安得拉邦学校委员会（今安得拉邦中等教育委员会）榜首。\(^\text{[14]}\)

1902年，拉曼进入马德拉斯（今钦奈）的总督学院就读，当时他的父亲已调任至该校教授数学与物理。\(^\text{[15]}\)1904年，他以优异成绩从马德拉斯大学获得文学学士（B.A.）学位，并在物理学与英语两门科目中均名列第一，荣获金牌奖。\(^\text{[14]}\)18岁时，仍为研究生的拉曼在英国期刊《哲学杂志》上发表了他的第一篇科学论文——《由矩形孔引起的不对称衍射带》（Unsymmetrical diffraction bands due to a rectangular aperture，1906年）。\(^\text{[16]}\)1907年，他以最高荣誉从同一所大学获得文学硕士（M.A.）学位。\(^\text{[17][18]}\)同年，他的第二篇论文发表于同一杂志，题为《液体的表面张力》。\(^\text{[19]}\)巧合的是，这篇论文与瑞利勋爵的一篇关于人耳对声音敏感度的论文同时刊登。\(^\text{[20]}\)由此，瑞利勋爵开始与拉曼通信，并礼貌地称呼他为“教授”。\(^\text{[14]}\)

拉曼的物理导师理查德·卢埃林·琼斯深知他的学术潜力，坚持建议他前往英国继续研究。\(^\text{[21]}\)琼斯为他安排了与吉法德上校的体检。然而，拉曼的健康状况欠佳，被认为是“体弱多病之人”。\(^\text{[22]}\)体检结果显示他无法适应英国严酷的气候。\(^\text{[11]}\)对此，拉曼后来回忆道：“吉法德检查了我，并证明如果我去英国，我会死于肺结核。”\(^\text{[23]}\)
\subsection{职业生涯}
拉曼的哥哥钱德拉塞卡拉·苏布拉马尼亚·艾耶早年已加入享有盛誉的印度政府部门——印度财政服务处（Indian Finance Service，今为印度审计与会计服务 Indian Audit and Accounts Service）。\(^\text{[24]}\)拉曼也步其后尘，于1907年2月的入职考试中名列第一，顺利取得任职资格。\(^\text{[25]}\)同年6月，他被派往加尔各答（今加尔各答）任副会计长。\(^\text{[11]}\)

在加尔各答期间，拉曼对印度科学促进会（Indian Association for the Cultivation of Science，简称 IACS）深感钦佩。\(^\text{[23]}\)该机构创立于1876年，是印度历史上第一所科学研究机构。他很快结识了几位重要人物：阿修托什·德伊，后来的长期合作伙伴；阿姆里塔·拉尔·瑟卡尔），IACS创办人马亨德拉拉尔·瑟卡尔之子、协会秘书；阿修托什·穆克吉，IACS执行委员、加尔各答大学校长。\(^\text{[14]}\)在他们的支持下，拉曼获准利用业余时间在IACS进行科学研究——哪怕是在“极其不寻常的时间”，\(^\text{[26]}\)正如他后来回忆的那样。\(^\text{[11]}\)当时，IACS尚未正式招募全职研究人员，也从未发表过科研论文。拉曼在1907年发表的论文《偏振光下的牛顿环》发表于《自然》杂志，成为该研究所的第一篇学术论文。\(^\text{[27]}\)这项成果极大地激发了IACS的科研热情，并促使其于1909年创办了《印度科学促进会公报》，其中拉曼是主要撰稿人。\(^\text{[14]}\)

1909年，拉曼被调任至英属缅甸仰光（今缅甸仰光），担任货币官。然而不久之后，他因父亲病逝而返回马德拉斯。父亲的去世及随后的丧礼使他在当地滞留了一整年。\(^\text{[28]}\)不久，他重新回到仰光任职，但在1910年又被调回印度马哈拉施特拉邦的那格浦尔。\(^\text{[29]}\)尚未在那格浦尔任职满一年，1911年他便被晋升为会计长，再次调往加尔各答。\(^\text{[28]}\)

从1915年起，加尔各答大学开始将研究生派至印度科学促进会（IACS），在拉曼的指导下进行科研工作。他的第一位学生是苏丹舒·库马尔·班纳吉（Sudhangsu Kumar Banerji），当时是加内什·普拉萨德指导下的博士研究生，后来成为印度气象局天文台总监。

自次年起，其他大学也纷纷效仿，如阿拉哈巴德大学、仰光大学、印多尔皇后学院、那格浦尔科学研究所、克里什纳斯学院以及马德拉斯大学等，\(^\text{[30]}\)都派学生前往IACS接受拉曼的指导。到1919年为止，拉曼已指导过十多名学生。\(^\text{[31]}\)同年，IACS秘书阿姆里塔·拉尔·瑟卡尔去世后，拉曼被授予两个荣誉职务：名誉教授与名誉秘书。\(^\text{[26]}\)拉曼将这一时期称为他人生的“黄金时代”。\(^\text{[32]}\)

加尔各答大学于1913年选定拉曼担任“帕利特物理学教授”，该职位由捐助人塔拉克纳特·帕利特爵士）设立。大学参议院于1914年1月30日正式任命，如会议记录所载：\\\\
“在1914年1月30日的参议院会议上，作出以下帕利特教授职位的任命：P. C. 雷博士与C. V. 拉曼先生。……每位教授的任命为长期制。教授须于年满六十岁时离任。”\(^\text{[14]}\)

在1914年之前，阿修托什·穆克吉曾邀请贾格迪什·钱德拉·玻色出任该职位，但玻色婉拒了邀请。\(^\text{[14]}\)作为第二人选，拉曼被任命为首任“帕利特物理学教授”，但由于第一次世界大战的爆发，他的上任被迫推迟。直到1917年，他加入加尔各答大学于1914年新设的拉贾巴扎尔理学院，才正式成为全职教授。\(^\text{[14]}\)拉曼在担任政府公务员十年后，勉强辞去职务——这一决定被称为一种“至高的牺牲”\(^\text{[26]}\)，因为教授的薪水几乎只有他原先收入的一半。不过，他的任职条件在加入大学时被明确写入报告，对他极为有利。报告中指出：\\\\
“C. V. 拉曼先生接受了塔拉克纳特·帕利特爵士物理学教授一职，条件是他无需被派往印度以外的地区……报告显示，C. V. 拉曼先生自1917年7月2日起正式就任帕利特物理学教授……拉曼先生已通知，他不必承担硕士或理学硕士课程的教学任务，以免影响其科研工作或指导高年级研究生进行研究。”\(^\text{[30]}\)

拉曼被任命为“帕利特物理学教授”一事，曾遭到加尔各答大学参议院部分成员的强烈反对，尤其是外国成员——理由是他既没有博士学位，也从未出国留学。作为回应，阿修托什·穆克吉特地安排加尔各答大学于1921年授予拉曼名誉理学博士学位，以示肯定。同年，他受邀前往牛津大学，在“大英帝国大学大会”上发表演讲。\(^\text{[34]}\)到那时，拉曼已在国际上享有盛誉，其接待方正是两位诺贝尔奖得主：J. J. 汤姆孙爵士与卢瑟福勋爵。\(^\text{[35]}\)1924年，拉曼当选为英国皇家学会院士后，穆克吉问他今后的目标，他笑答道：“当然是诺贝尔奖。”\(^\text{[26]}\)1926年，拉曼创办了《印度物理学杂志》，并担任首任主编。\(^\text{[36]}\)该杂志第二卷刊登了他著名的论文《一种新的辐射》，首次报道了“拉曼效应”的发现。\(^\text{[37][38]}\)

1932年，拉曼的“帕利特物理学教授”职位由德本德拉·莫汉·玻色接任。翌年，他被任命为位于班加罗尔的印度科学研究所（Indian Institute of Science，简称 IISc）所长，于1933年离开加尔各答。\(^\text{[39]}\)印度科学研究所的创立得到了多方支持：迈索尔国王克里希纳拉贾·瓦迪亚尔四世、实业家贾姆谢特吉·塔塔以及海得拉巴的尼扎姆——米尔·奥斯曼·阿里·汗爵士共同捐赠了土地与资金。1909年，印度总督明托勋爵批准建立该研究所，并由英国政府任命莫里斯·特拉弗斯为首任所长。\(^\text{[40]}\)拉曼成为第四任所长，也是第一位印度籍所长。在任期间，他聘用了G. N. 拉马钱德兰，后者后来成为著名的X射线晶体学家。1934年，拉曼创立了印度科学院，并创办了其学术期刊《印度科学院院刊》，该期刊后来分为三个分支刊物：《数学科学院刊》，《化学学报》，《地球系统科学学报》。\(^\text{[35]}\)大约在同一时期，拉曼早在1917年就提出设想的“加尔各答物理学会”也正式成立。\(^\text{[14]}\)

1943年，拉曼与他的前学生潘查帕凯萨·克里希纳穆尔蒂共同创办了“特拉凡哥尔化学制造有限公司”。\(^\text{[41]}\)该公司于1996年更名为“TCM有限公司”，是印度最早的有机与无机化学品制造企业之一。\(^\text{[42]}\)1947年，印度独立后新政府任命拉曼为首任“国家教授”。\(^\text{[43]}\)

拉曼于1948年从印度科学研究所退休，次年在班加罗尔创立了“拉曼研究所”。他担任所长，并一直在此工作，直至1970年去世。\(^\text{[43]}\)
\subsection{科学贡献}
\subsubsection{音乐声学}
\begin{figure}[ht]
\centering
\includegraphics[width=8cm]{./figures/3a7f4c6530b8aaba.png}
\caption{能级示意图：展示了参与拉曼信号的各个能级状态} \label{fig_QDLlm_1}
\end{figure}
拉曼的研究兴趣之一是音乐声音的科学基础。他受到赫尔曼·冯·亥姆霍兹所著《音的感觉》一书的启发，这本书是他加入印度科学促进会（IACS）后接触到的。\(^\text{[25]}\)在1916年至1921年间，拉曼大量发表了相关研究成果。他提出了基于速度叠加原理的弓弦乐器横向振动理论。他最早的研究之一是关于小提琴和大提琴中的“狼音”现象。\(^\text{[44][45]}\)此外，他还研究了各种小提琴及类似乐器的声学特性，包括印度传统弦乐器的声学行为。\(^\text{[46][47]}\)他甚至研究了水滴溅起的声音。\(^\text{[48]}\)拉曼还进行了他称之为“机械演奏小提琴实验”的研究。\(^\text{[49]}\)
\begin{figure}[ht]
\centering
\includegraphics[width=8cm]{./figures/8cb85b7f2b42ce80.png}
\caption{1930年诺贝尔奖颁奖典礼上的拉曼（左起第一位），与其他获奖者合影：C. V. 拉曼（物理学奖）、汉斯·费舍尔（化学奖）、卡尔·兰德施泰纳（医学奖）以及辛克莱·刘易斯（文学奖）。} \label{fig_QDLlm_2}
\end{figure}
拉曼还研究了印度鼓类乐器的独特性。\(^\text{[50]}\)他对塔布拉鼓与姆里当加姆鼓声音谐波特性的分析，是印度打击乐器声学的首批科学研究。\(^\text{[51]}\)他还撰写了一篇关于钢琴琴弦振动的批判性论文，探讨并修正了被称为考夫曼理论的观点。\(^\text{[52]}\)1921年短暂访问英国期间，拉曼抓住机会研究了伦敦圣保罗大教堂圆顶“回音廊”中的声音传播现象，该建筑能产生奇特的回声与声学效果。\(^\text{[53][54]}\)拉曼在声学方面的研究，无论在实验方法还是概念上，都为他后来的光学与量子力学研究奠定了重要基础。\(^\text{[55]}\)
\subsubsection{海水的蓝色}
在拓展光学研究的过程中，拉曼自1919年起开始研究光的散射现象。\(^\text{[56]}\)他在光学物理方面的第一个重要发现，便是对海水蓝色成因的解释。1921年9月，他从英国乘坐S.S. Narkunda号轮船返回印度途中，凝视地中海时，对海水的蓝色产生了浓厚兴趣。凭借随身携带的简单光学仪器——一台袖珍分光镜和一个尼科尔棱镜，他亲自对海水进行了观察与研究。\(^\text{[57]}\)当时，关于海水为何呈蓝色的多种假说中，\(^\text{[58][59]}\)最广为接受的是瑞利勋爵在1910年提出的解释：“深海那令人赞叹的深蓝色，与水本身的颜色无关，而只是天空蓝光在海面的反射。”\(^\text{[60]}\)瑞利对天空为何呈蓝色的解释是正确的——这一现象即如今所称的“瑞利散射”，即光在大气中被微粒散射与折射所致。\(^\text{[61][62]}\)因此，他关于海水蓝色来源于天空反射的观点在当时被普遍认为是正确的。然而，拉曼通过使用尼科尔棱镜消除了海面反射阳光的影响，直接观察海水后发现：海水在这种条件下显得比平时更加湛蓝。这一结果与瑞利的理论相矛盾，也使他确信，海水的蓝色并非仅仅来自天空的反射。\(^\text{[63]}\)

当 *S.S. Narkunda* 号轮船抵达孟买港（今孟买港）后，拉曼立即完成了一篇题为《海的颜色》（*The Colour of the Sea*）的文章，并发表在1921年11月的《自然》（*Nature*）杂志上。他在文中指出，瑞利的解释“通过一种简单的观察方式（使用尼科尔棱镜）就可以被质疑”。【63】他写道：\\\\
“当通过一个放置在眼前的尼科尔棱镜观察海水，以消除表面反射时，可以看到阳光的光线穿入水中，并由于透视效应似乎在水的深处汇聚成一点。问题是：究竟是什么在散射光线，使得它的路径变得可见？一个值得考虑的有趣可能性是——这些散射光的粒子，至少部分，可能正是水分子本身。”【14】
\begin{figure}[ht]
\centering
\includegraphics[width=6cm]{./figures/c373caa8734590ef.png}
\caption{} \label{fig_QDLlm_3}
\end{figure}